% Part: first-order-logic
% Chapter: proof-systems
% Section: introduction

\documentclass[../../../include/open-logic-section]{subfiles}

\begin{document}

\iftag{FOL}
      {\olfileid[hi]{fol}{prf}{int}}
      {\olfileid[hi]{pl}{prf}{int}}

\olsection{परिचय}

तर्कशास्त्रों में सामान्यतः अर्थविज्ञान और !!{derivation}-तंत्र दोनों होते
हैं। अर्थविज्ञान सत्य, संतुष्टनीयता, वैधता और तार्किक निष्पत्ति जैसी धारणाओं
से संबंधित है। !!{derivation}-तंत्रों का उद्देश्य तार्किक निष्पत्ति और वैधता
स्थापित करने की पूर्णतः वाक्यविन्यासगत विधि देना है। वे इस अर्थ में पूर्णतः
वाक्यविन्यासगत हैं कि ऐसे तंत्र की !!{derivation} एक परिमित वाक्यविन्यासगत
वस्तु होती है, प्रायः !!{sentence}s या !!{formula}s का अनुक्रम (या कोई अन्य
परिमित विन्यास)। अच्छे !!{derivation}-तंत्रों में यह गुण होता है कि
!!{sentence}s या !!{formula}s के किसी भी दिए अनुक्रम अथवा विन्यास के ``सही''
होने को यांत्रिक रूप से जाँचा जा सकता है।

प्रथम-कोटि तर्कशास्त्र के सबसे सरल (और ऐतिहासिक रूप से प्रथम)
!!{derivation}-तंत्र \emph{अभिगृहीतीय} थे। ऐसे तंत्र में !!{formula}s का कोई
अनुक्रम तभी !!{derivation} माना जाता है, जब उसमें हर अलग !!{formula} या तो
``अभिगृहीतों'' के किसी नियत समुच्चय में हो, अथवा अनुक्रम में उससे पहले आए
!!{formula}s से नियत संख्या में से किसी ``अनुमान-नियम'' द्वारा निकलता हो—और
यह यांत्रिक रूप से जाँचा जा सकता है कि !!{formula} अभिगृहीत है या नहीं तथा
वह किसी अनुमान-नियम द्वारा अन्य !!{formula}s से ठीक प्रकार निकलता है या नहीं।
अभिगृहीतीय !!{derivation}-तंत्रों का वर्णन करना आसान है—और अधिसैद्धांतिक रूप
से उन्हें सँभालना भी आसान है—लेकिन उनमें !!{derivation}s पढ़ना और समझना कठिन
है तथा उन्हें बनाना भी कठिन है।

अन्य !!{derivation}-तंत्र इसलिए विकसित किए गए हैं कि !!{derivation}s बनाना
आसान हो या पूरी हो चुकी !!{derivation}s अधिक आसानी से समझ में आएँ। उदाहरण हैं
स्वाभाविक निगमन, सत्य-वृक्ष—जिन्हें सारणी-प्रमाण भी कहते हैं—और अनुक्रम-कलन।
कुछ !!{derivation}-तंत्र विशेषतः यंत्रीकरण को ध्यान में रखकर बनाए जाते हैं;
जैसे, खंड-समाधान विधि को संगणक-कार्यक्रम में लागू करना आसान है (लेकिन उसकी
!!{derivation}s को समझना लगभग असंभव है)। इन अन्य !!{derivation}-तंत्रों में से
अधिकतर में !!{derivation}s, !!{formula}s के अनुक्रमों के बजाय, वृक्षों के रूप
में प्रस्तुत होती हैं। इससे यह देखना आसान होता है कि !!{derivation} के
कौन-से भाग किन दूसरे भागों पर निर्भर हैं।

अतः किसी दिए तर्कशास्त्र—जैसे प्रथम-कोटि तर्कशास्त्र—के अलग-अलग
!!{derivation}-तंत्र इस बात की अलग-अलग स्पष्ट व्याख्याएँ देंगे कि
!!{sentence} का \emph{प्रमेय} होना क्या है और किसी !!{sentence} का कुछ
अन्य वाक्यों से !!{derivable} होना क्या है। यह चाहे जिस तरह किया जाए
(अभिगृहीतीय !!{derivation}s, स्वाभाविक निगमन, अनुक्रम !!{derivation}s,
सत्य-वृक्ष या खंड-समाधान द्वारा खंडन), हम चाहते हैं कि ये संबंध वैधता और
तार्किक निष्पत्ति की अर्थवैज्ञानिक धारणाओं से मेल खाएँ। $\Proves !A$ को
``$!A$ एक प्रमेय है'' के लिए और ``$\Gamma \Proves !A$'' को ``$!A$,
$\Gamma$ से !!{derivable} है'' के लिए लिखें। $\Proves$ को चाहे जिस प्रकार
परिभाषित करें, हम चाहते हैं कि यह $\Entails$ से मेल खाए; अर्थात्:
\begin{enumerate}
\item $\Proves !A$ तभी और केवल तभी जब $\Entails !A$
\item $\Gamma \Proves !A$ तभी और केवल तभी जब $\Gamma \Entails !A$
\end{enumerate}
ऊपर दी गई ``केवल तभी'' दिशा को \emph{सुदृढ़ता} कहते हैं।
!!^{derivation}-तंत्र सुदृढ़ है यदि !!{derivability} तार्किक निष्पत्ति (या
वैधता) की गारंटी देती है। हर उपयोगी !!{derivation}-तंत्र को सुदृढ़ होना चाहिए;
असुदृढ़ !!{derivation}-तंत्र बिल्कुल उपयोगी नहीं होते। आखिर !!{derivation}
का पूरा उद्देश्य वैधता या तार्किक निष्पत्ति की वाक्यविन्यासगत गारंटी देना है।
हम प्रस्तुत किए गए !!{derivation}-तंत्रों की सुदृढ़ता सिद्ध करेंगे।

इसकी विपरीत ``यदि'' दिशा भी महत्त्वपूर्ण है; इसे \emph{पूर्णता} कहते हैं। कोई
पूर्ण !!{derivation}-तंत्र इतना शक्तिशाली होता है कि जब भी $!A$ वैध हो तब
$!A$ के प्रमेय होने को, और $\Gamma \Proves !A$ को, जब भी
$\Gamma \Entails !A$, दिखा सके। पूर्णता स्थापित करना अधिक कठिन है, और कुछ
तर्कशास्त्रों के कोई पूर्ण !!{derivation}-तंत्र नहीं होते। प्रथम-कोटि
तर्कशास्त्र का तंत्र पूर्ण है। कुर्ट गेडेल ने अपने 1929 के शोधप्रबंध में
प्रथम-कोटि तर्कशास्त्र के !!{derivation}-तंत्र की पूर्णता सबसे पहले सिद्ध की।

!!{derivation}-तंत्रों से जुड़ी एक अन्य धारणा \emph{संगतता} है।
!!{sentence}s का कोई समुच्चय असंगत कहलाता है यदि उससे कुछ भी !!{derive} किया
जा सकता हो; अन्यथा वह संगत है। असंगतता, असंतुष्टनीयता का
\mbox{वाक्यविन्यासीय}
प्रतिरूप है: असंतुष्टनीय समुच्चयों की तरह, !!{sentence}s के असंगत समुच्चय
अच्छे सिद्धान्त नहीं बनते; उनमें मूलभूत दोष होता है। !!{sentence}s के संगत
समुच्चय सत्य या उपयोगी न भी हों, फिर भी वे तार्किक उपयोगिता की न्यूनतम कसौटी
पूरी करते हैं। अलग-अलग !!{derivation}-तंत्रों में !!{sentence}s के
समुच्चयों की संगतता की विशिष्ट परिभाषा अलग हो सकती है; लेकिन~$\Proves$ की तरह
हम चाहते हैं कि संगतता अपने अर्थवैज्ञानिक प्रतिरूप, संतुष्टनीयता, से मिले। हम
चाहते हैं कि सदा $\Gamma$ संगत तभी और केवल तभी हो जब वह संतुष्टनीय हो। यहाँ
``केवल तभी'' दिशा पूर्णता के बराबर है (संगतता संतुष्टनीयता की गारंटी देती है),
और ``यदि'' दिशा सुदृढ़ता के बराबर है (संतुष्टनीयता संगतता की गारंटी देती है)।
वास्तव में, शास्त्रीय प्रथम-कोटि तर्कशास्त्र में सुदृढ़ता और पूर्णता के दोनों
रूप समतुल्य हैं।

\end{document}
